\clearpage
\phantomsection% 
\addcontentsline{toc}{chapter}{ABSTRACT}
\thispagestyle{fancy}

\begin{center}
	\large \bfseries \MakeUppercase{Abstract}\\
	\normalsize \normalfont {\thetitleEN}\\
	\normalsize \normalfont {\theauthor}\\
	\bigskip
	
	\normalsize \normalfont \justifying \singlespacing
	The advancement of astronomy observation has shifted toward automation systems, yet many automated telescopes face challenges in initial configuration and calibration, particularly for users with limited technical skills. This research develops a robotic telescope based on Alt-Azimuth mount with automatic control system through an Android application as the main interface. The system integrates IMU sensor MPU6050, magnetometer QMC5883L, and magnetic encoder AS5600 to provide real-time orientation information, enabling the telescope to automatically adjust its position when powered on or shifted. To address the limitations of MPU6050 sensor (noise, drift, and inaccuracy), Mahony Filter is implemented, which combines gyroscope and accelerometer data using a PI (Proportional-Integral) controller approach to produce more accurate orientation estimates. System performance is evaluated using standard deviation metrics to measure sensor reading stability. The results show that Mahony Filter successfully reduces Roll noise by 79.1\% (from 0.510$^{\circ}$ to 0.107$^{\circ}$) and Pitch noise by 84.2\% (from 1.719$^{\circ}$ to 0.271$^{\circ}$). The GoTo system successfully directs the telescope to celestial objects with pointing errors ranging from 0$^{\circ}$10' to 2$^{\circ}$30', while the tracking system can follow the apparent motion of celestial objects with varied results. Solar observation for 19 minutes showed the object remained within the field of view, however, tracking objects near zenith experienced significant drift causing objects to move out of the field of view. These limitations are caused by a combination of mechanical factors (PLA structure flex and gear backlash), limited update frequency (5-10 seconds), and error accumulation. This system is adequate for amateur astronomy and education applications with bright celestial objects of large angular diameter at positions not near zenith, but still requires significant mechanical improvements for observing point objects (planets, stars) that require higher precision.\par
    
    \vspace{20pt}
	\raggedright\textbf{Keywords: Robotic Telescope, Alt-Azimuth Mount, Mahony Filter, IMU Sensor, GoTo System}
	
	\vfill
	
\end{center}
\clearpage